\chapter{2014年福建卷（理）}

\section{单选题}

\begin{problem}
复数 \(z = (3 - 2\i)\i\) 的共轭复数 \(\overline{z}\) 等于（\quad）

\choices
    {\(-2 - 3\i\)}
    {\(-2 + 3\i\)}
    {\(2 - 3\i\)}
    {\(2 + 3\i\)}
\end{problem}

\begin{answer}
C
\end{answer}

\begin{solution}
由复数乘法，\(z=(3-2\i)\i=3\i-2\i^2=2+3\i\)，所以 \(\overline{z}=2-3\i\). 故选 C.
\end{solution}

\begin{problem}
某空间几何体的正视图是三角形，则该几何体不可能是（\quad）

\choices
    {圆柱}
    {圆锥}
    {四面体}
    {三棱柱}
\end{problem}

\begin{answer}
A
\end{answer}

\begin{solution}
圆柱的正视图为矩形，不能为三角形；圆锥的正视图可以为三角形，四面体的正视图也可以为三角形，三棱柱的正视图还可以为三角形. 因此不可能的是圆柱，故选 A.
\end{solution}

\begin{problem}
等差数列 \(\{a_{n}\}\) 的前 \(n\) 项和为 \(S_{n}\)，若 \(a_{1} = 2\)，\(S_{3} = 12\)，则 \(a_{6} =\)（\quad）

\choices
    {8}
    {10}
    {12}
    {14}
\end{problem}

\begin{answer}
C
\end{answer}

\begin{solution}
设等差数列的公差为 \(d\). 由前 \(3\) 项和公式，得 \(S_3=3a_1+3d=3\times2+3d=12\)，所以 \(d=2\). 因而 \(a_6=a_1+5d=2+5\times2=12\)，故选 C.
\end{solution}

\begin{problem}
若函数 \(y = \log_a x\)（\(a > 0\)，且 \(a \neq 1\)）的图象如图所示，则下列函数图象正确的是（\quad）

\bitmapfigure[width=0.42\linewidth]{img/2014/fujian_science/q4_fig1.png}

\choices
    {\choicebitmap[width=0.15\paperwidth]{img/2014/fujian_science/q4_optA.png}}
    {\choicebitmap[width=0.15\paperwidth]{img/2014/fujian_science/q4_optB.png}}
    {\choicebitmap[width=0.15\paperwidth]{img/2014/fujian_science/q4_optC.png}}
    {\choicebitmap[width=0.15\paperwidth]{img/2014/fujian_science/q4_optD.png}}
\end{problem}

\begin{answer}
B
\end{answer}

\begin{solution}
由题图，函数 \(y=\log_a x\) 经过点 \((3,1)\)，所以 \(\log_a3=1\)，从而 \(a=3\). 因此选项 A 中的 \(y=a^{-x}=3^{-x}\) 应为减函数，图象不符；选项 B 中的 \(y=x^a=x^3\) 经过 \((0,0)\)、\((1,1)\)，且在整个定义域上递增，图象与选项 B 相符；选项 C 的 \(y=(-x)^a=(-x)^3\) 应为减函数，图象不符；选项 D 的定义域应为 \(x<0\)，且当 \(x\to0^-\) 时 \(\log_a(-x)\to-\infty\)，图象不符. 故选 B.
\end{solution}

\begin{problem}
阅读如图所示的程序框图，运行相应的程序，输出的 \(S\) 的值等于（\quad）

\texfigure[max width=0.42\linewidth]{img_repaint/2014/fujian_science/q5_fig1.tex}

\choices
    {18}
    {20}
    {21}
    {40}
\end{problem}

\begin{answer}
B
\end{answer}

\begin{solution}
程序开始时 \(S=0\)，\(n=1\). 第一次循环后 \(S=0+2^1+1=3\)，\(n=2\)；第二次循环后 \(S=3+2^2+2=9\)，\(n=3\)；第三次循环后 \(S=9+2^3+3=20\)，\(n=4\). 此时 \(S\geqslant15\)，程序输出 \(S=20\). 故选 B.
\end{solution}

\begin{problem}
直线 \(l: y = kx + 1\) 与圆 \(O: x^2 + y^2 = 1\) 相交于 \(A\)，\(B\) 两点，则“\(k = 1\)”是“\(\triangle OAB\) 的面积为 \(\frac{1}{2}\)”的（\quad）

\choices
    {充分而不必要条件}
    {必要而不充分条件}
    {充分必要条件}
    {既不充分又不必要条件}
\end{problem}

\begin{answer}
A
\end{answer}

\begin{solution}
直线 \(l\) 的方程可写为 \(kx-y+1=0\)，所以原点到直线 \(l\) 的距离为 \(h=\frac{1}{\sqrt{k^2+1}}\). 由弦心距公式，弦 \(AB\) 的长度为
\(AB=2\sqrt{1-h^2}=2\sqrt{1-\frac{1}{k^2+1}}=\frac{2|k|}{\sqrt{k^2+1}}\).

当直线与圆有两个交点时，\(k\neq0\)，于是
\(S_{\triangle OAB}=\frac{1}{2}AB\cdot h=\frac{1}{2}\cdot\frac{2|k|}{\sqrt{k^2+1}}\cdot\frac{1}{\sqrt{k^2+1}}=\frac{|k|}{k^2+1}\).
若 \(k=1\)，则 \(S_{\triangle OAB}=\frac{1}{2}\)，所以“\(k=1\)”是该面积条件的充分条件. 反之，若面积为 \(\frac{1}{2}\)，则 \(\frac{|k|}{k^2+1}=\frac{1}{2}\)，即 \(k^2-2|k|+1=0\)，从而 \(|k|=1\)，既可能有 \(k=1\)，也可能有 \(k=-1\). 因此该条件不是必要条件，故选 A.
\end{solution}

\begin{problem}
已知函数 \( f(x) = \begin{cases}
x^2 + 1, & x > 0,\\
\cos x, & x \leqslant 0.
\end{cases} \) 则下列结论正确的是（\quad）

\choices
    {\(f(x)\) 是偶函数}
    {\(f(x)\) 是增函数}
    {\(f(x)\) 是周期函数}
    {\(f(x)\) 的值域为 \([-1, +\infty)\)}
\end{problem}

\begin{answer}
D
\end{answer}

\begin{solution}
当 \(x>0\) 时，\(f(x)=x^2+1>1\)，其值域为 \((1,+\infty)\)；当 \(x\leqslant0\) 时，\(f(x)=\cos x\)，其值域为 \([-1,1]\). 所以 \(f(x)\) 的值域为 \([-1,+\infty)\)，故选 D.
\end{solution}

\begin{problem}
在下列向量组中，可以把向量 \(\bs{a} = (3, 2)\) 表示出来的是（\quad）

\choices
    {\(\bs{e}_1 = (0,0), \bs{e}_2 = (1,2)\)}
    {\(\bs{e}_1 = (-1,2)\)，\(\bs{e}_2 = (5, -2)\)}
    {\(\bs{e}_1 = (3,5)\)，\(\bs{e}_2 = (6,10)\)}
    {\(\bs{e}_1 = (2, -3)\)，\(\bs{e}_2 = (-2, 3)\)}
\end{problem}

\begin{answer}
B
\end{answer}

\begin{solution}
选项 A 中 \(\bs{e}_1=(0,0)\)，任意线性组合都形如 \(\lambda(1,2)\). 若其等于 \(\bs{a}=(3,2)\)，由第一坐标得 \(\lambda=3\)，但第二坐标应为 \(6\)，矛盾，所以选项 A 不能表示 \(\bs{a}\). 选项 C 中 \(\bs{e}_2=2\bs{e}_1\)，任意线性组合都形如 \(\lambda(3,5)\)；由第一坐标得 \(\lambda=1\) 时第二坐标为 \(5\)，所以也不能表示 \(\bs{a}\). 选项 D 中 \(\bs{e}_2=-\bs{e}_1\)，任意线性组合都形如 \(\lambda(2,-3)\)；由第一坐标得 \(\lambda=\frac{3}{2}\) 时第二坐标为 \(-\frac{9}{2}\)，同样不能表示 \(\bs{a}\).

选项 B 中，两个向量的行列式为 \(\begin{vmatrix}-1&5\\2&-2\end{vmatrix}=(-1)(-2)-5\times2=-8\neq0\)，所以它们不共线. 事实上，\(\bs{a}=2\bs{e}_1+\bs{e}_2=2(-1,2)+(5,-2)=(3,2)\)，故选 B.
\end{solution}

\begin{problem}
设 \(P\)，\(Q\) 分别为圆 \(x^{2} + (y - 6)^{2} = 2\) 和椭圆 \(\frac{x^2}{10} + y^2 = 1\) 上的点，则 \(P\)，\(Q\) 两点间的最大距离是（\quad）

\choices
    {\(5\sqrt{2}\)}
    {\(\sqrt{46} + \sqrt{2}\)}
    {\(7 + \sqrt{2}\)}
    {\(6\sqrt{2}\)}
\end{problem}

\begin{answer}
D
\end{answer}

\begin{solution}
设圆的圆心为 \(O_1=(0,6)\)，半径为 \(\sqrt{2}\). 对椭圆上的任意点 \(Q\)，圆上距 \(Q\) 最远的点 \(P\) 满足 \(PQ=QO_1+\sqrt{2}\)，因此只需先求 \(QO_1\) 的最大值.

设 \(Q=(\sqrt{10}\cos t,\sin t)\)，则
\(QO_1^2=10\cos^2t+(\sin t-6)^2=46-9\sin^2t-12\sin t\).
令 \(u=\sin t\)，则 \(-1\leqslant u\leqslant1\)，且
\(QO_1^2=46-9u^2-12u=50-9\left(u+\frac{2}{3}\right)^2\leqslant50\).
当 \(u=-\frac{2}{3}\) 时取等号，所以 \(QO_1\) 的最大值为 \(5\sqrt{2}\). 因而 \(PQ\) 的最大值为 \(5\sqrt{2}+\sqrt{2}=6\sqrt{2}\)，故选 D.
\end{solution}

\begin{problem}
用 \(a\) 代表红球，\(b\) 代表蓝球，\(c\) 代表黑球，由加法原理及乘法原理，从 \(1\) 个红球和 \(1\) 个蓝球中取出若干个球的所有取法可由 \((1 + a)(1 + b)\) 的展开式 \(1 + a + b + ab\) 表示出来. 如：“1”表示一个球都不取，“\(a\)”表示取出一个红球，而“\(ab\)”则表示把红球和蓝球都取出来. 以此类推，下列各式中，其展开式可用来表示从 \(5\) 个无区别的红球，\(5\) 个无区别的蓝球，\(5\) 个有区别的黑球中取出若干个球，且所有的蓝球都取出或都不取出的所有取法的是（\quad）

\choices
    {\((1 + a + a^2 + a^3 + a^4 + a^5)(1 + b^5)(1 + c)^5\)}
    {\((1 + a^5)(1 + b + b^2 + b^3 + b^4 + b^5)(1 + c)^5\)}
    {\((1 + a)^5(1 + b + b^2 + b^3 + b^4 + b^5)(1 + c^5)\)}
    {\((1 + a^5)(1 + b)^5(1 + c + c^2 + c^3 + c^4 + c^5)\)}
\end{problem}

\begin{answer}
A
\end{answer}

\begin{solution}
无区别的红球可以取出 \(0,1,2,\ldots,5\) 个，对应因子为 \(1+a+a^2+a^3+a^4+a^5\)；题目要求蓝球全部取出或全部不取出，对应因子为 \(1+b^5\)；\(5\) 个有区别的黑球中每个球都有取或不取两种选择，对应因子为 \((1+c)^5\). 因此所求展开式为
\((1+a+a^2+a^3+a^4+a^5)(1+b^5)(1+c)^5\)，故选 A.
\end{solution}

\section{填空题}

\begin{problem}
若变量 \(x\)，\(y\) 满足约束条件 \(\begin{cases}
x - y + 1 \leqslant 0,\\
x + 2y - 8 \leqslant 0,\\
x \geqslant 0,
\end{cases}\) 则 \(z = 3x + y\) 的最小值为\fillinblank{}.
\end{problem}

\begin{answer}
1
\end{answer}

\begin{solution}
由约束条件，得 \(y\geqslant x+1\)，\(y\leqslant4-\frac{x}{2}\)，且 \(x\geqslant0\). 因而可行域的顶点为 \((0,1)\)、\((0,4)\)、\((2,3)\). 在这三个顶点处，目标函数 \(z=3x+y\) 的值分别为 \(1\)、\(4\)、\(9\). 所以 \(z\) 的最小值为 \(1\)，故填 \(1\).
\end{solution}

\begin{problem}
在 \(\triangle ABC\) 中，\(A = 60^{\circ}\)，\(AC = 4\)，\(BC = 2\sqrt{3}\)，则 \(\triangle ABC\) 的面积等于\fillinblank{}.
\end{problem}

\begin{answer}
\(2\sqrt{3}\)
\end{answer}

\begin{solution}
设 \(AB=c\). 由余弦定理，得
\(BC^2=AB^2+AC^2-2AB\cdot AC\cos A\)，即 \((2\sqrt{3})^2=c^2+4^2-2\cdot c\cdot4\cdot\cos60^{\circ}\). 整理得 \(12=c^2+16-4c\)，所以 \((c-2)^2=0\)，从而 \(AB=2\).

因此 \(\triangle ABC\) 的面积为 \(\frac{1}{2}AB\cdot AC\sin A=\frac{1}{2}\times2\times4\times\frac{\sqrt{3}}{2}=2\sqrt{3}\)，故填 \(2\sqrt{3}\).
\end{solution}

\begin{problem}
要制作一个容积为 \(4 \mathrm{~m}^{3}\)，高为 \(1 \mathrm{~m}\) 的无盖长方体容器，已知该容器的底面造价是每平方米 \(20\text{ 元}\)，侧面造价是每平方米 \(10\text{ 元}\)，则该容器的最低总造价是\fillinblank{}\(\text{元}\).
\end{problem}

\begin{answer}
160
\end{answer}

\begin{solution}
设长方体底面的两条边长分别为 \(x\mathrm{~m}\)、\(y\mathrm{~m}\)，则 \(x>0\)，\(y>0\)，且由容积条件得 \(xy\times1=4\)，即 \(xy=4\).

底面造价为 \(20xy=80\) 元，四个侧面的总面积为 \(2x\times1+2y\times1=2(x+y)\)，侧面造价为 \(20(x+y)\) 元. 所以总造价为 \(C=80+20(x+y)\). 由基本不等式，\(x+y\geqslant2\sqrt{xy}=4\)，当且仅当 \(x=y=2\) 时取等号. 因而总造价的最小值为 \(80+20\times4=160\) 元，故填 \(160\).
\end{solution}

\begin{problem}
如图，在边长为 \(\e\)（\(\e\) 为自然对数的底数）的正方形中随机撒一粒黄豆，则它落到阴影部分的概率为\fillinblank{}.

\bitmapfigure[max width=0.30\linewidth]{img/2014/fujian_science/q14_fig1.png}
\end{problem}

\begin{answer}
\(\frac{2}{\e^2}\)
\end{answer}

\begin{solution}
正方形边长为 \(\e\)，面积为 \(\e^2\). 由图，左上方阴影部分的面积为
\(\displaystyle\int_0^1(\e-\e^x)\,\mathrm{d}x=\e-(\e-1)=1\)，右下方阴影部分的面积为
\(\displaystyle\int_1^{\e}\ln x\,\mathrm{d}x=\left[x\ln x-x\right]_1^{\e}=1\).
因此阴影部分总面积为 \(2\)，黄豆落在阴影部分的概率为 \(\frac{2}{\e^2}\)，故填 \(\frac{2}{\e^2}\).
\end{solution}

\begin{problem}
若集合 \(\{a, b, c, d\} = \{1, 2, 3, 4\}\)，且下列四个关系：

\begin{circlelist}
\item \(a = 1\)；
\item \(b \neq 1\)；
\item \(c = 2\)；
\item \(d \neq 4\).
\end{circlelist}

有且只有一个是正确的，则符合条件的有序数组 \((a, b, c, d)\) 的个数是\fillinblank{}.
\end{problem}

\begin{answer}
6
\end{answer}

\begin{solution}
若 \(a=1\)，由于 \(a,b,c,d\) 是 \(1,2,3,4\) 的排列，必有 \(b\neq1\)，此时前两个关系同时正确，不符合“有且只有一个正确”. 所以必须有 \(a\neq1\).

若只有第（2）个关系正确，则 \(b\neq1\)，\(c\neq2\)，且第（4）个关系错误，即 \(d=4\). 在 \(a,b,c\) 中排列 \(1,2,3\)，满足 \(a\neq1\)、\(b\neq1\)、\(c\neq2\) 的取法为 \((a,b,c)=(2,3,1)\) 或 \((3,2,1)\)，共 \(2\) 个.

若只有第（3）个关系正确，则 \(b=1\)，\(c=2\)，且 \(d=4\)，从而 \(a=3\)，共 \(1\) 个.

若只有第（4）个关系正确，则 \(b=1\)，\(c\neq2\)，\(d\neq4\). 在 \(a,c,d\) 中排列 \(2,3,4\)：当 \(c=3\) 时只能取 \((a,d)=(4,2)\)；当 \(c=4\) 时 \((a,d)\) 可为 \((2,3)\) 或 \((3,2)\)，共 \(3\) 个.

所以符合条件的有序数组共有 \(2+1+3=6\) 个，故填 \(6\).
\end{solution}

\section{解答题}

\begin{problem}
已知函数 \( f(x) = \cos x (\sin x + \cos x) - \frac{1}{2} \)，

\begin{enumerate}
\item 若 \(0 < \alpha < \frac{\pi}{2}\)，且 \(\sin \alpha = \frac{\sqrt{2}}{2}\)，求 \(f(\alpha)\) 的值；
\item 求函数 \(f(x)\) 的最小正周期及单调递增区间.
\end{enumerate}
\end{problem}

\begin{solution}
\begin{enumerate}
\item 因为 \(0<\alpha<\frac{\pi}{2}\)，且 \(\sin\alpha=\frac{\sqrt{2}}{2}\)，所以 \(\alpha=\frac{\pi}{4}\)，从而 \(\cos\alpha=\frac{\sqrt{2}}{2}\). 因此
\(f(\alpha)=\cos\alpha(\sin\alpha+\cos\alpha)-\frac{1}{2}=\frac{\sqrt{2}}{2}\left(\frac{\sqrt{2}}{2}+\frac{\sqrt{2}}{2}\right)-\frac{1}{2}=\frac{1}{2}\).

\item 利用二倍角公式，得
\(f(x)=\sin x\cos x+\cos^2x-\frac{1}{2}=\frac{1}{2}\sin2x+\frac{1+\cos2x}{2}-\frac{1}{2}=\frac{1}{2}(\sin2x+\cos2x)=\frac{\sqrt{2}}{2}\sin\left(2x+\frac{\pi}{4}\right)\).

因此函数的最小正周期为 \(T=\frac{2\pi}{2}=\pi\). 又
\(f'(x)=\cos2x-\sin2x=\sqrt{2}\cos\left(2x+\frac{\pi}{4}\right)\).
当 \(k\in\Z\) 时，\(f'(x)>0\) 等价于
\(-\frac{\pi}{2}+2k\pi<2x+\frac{\pi}{4}<\frac{\pi}{2}+2k\pi\)，即
\(k\pi-\frac{3\pi}{8}<x<k\pi+\frac{\pi}{8}\). 所以 \(f(x)\) 的单调递增区间为 \(\left(k\pi-\frac{3\pi}{8},k\pi+\frac{\pi}{8}\right)\)，\(k\in\Z\).
\end{enumerate}
\end{solution}

\begin{problem}
在平面四边形 \(ABCD\) 中，\(AB = BD = CD = 1\)，\(AB \perp BD\)，\(CD \perp BD\). 将 \(\triangle ABD\) 沿 \(BD\) 折起，使得平面 \(ABD \perp\) 平面 \(BCD\)，如图.

\begin{enumerate}
\item 求证：\(AB \perp CD\)；
\item 若 \(M\) 为 \(AD\) 中点，求直线 \(AD\) 与平面 \(MBC\) 所成角的正弦值.
\end{enumerate}

\bitmapfigure[max width=0.35\linewidth]{img/2014/fujian_science/q17_fig1.png}
\end{problem}

\begin{solution}
\begin{enumerate}
\item 因为平面 \(ABD\perp\) 平面 \(BCD\)，且两平面的交线为 \(BD\)，又 \(CD\) 在平面 \(BCD\) 内且 \(CD\perp BD\)，由二面角的性质可知 \(CD\perp\) 平面 \(ABD\). 由于 \(AB\) 在平面 \(ABD\) 内，所以 \(CD\perp AB\).

\item 建立空间直角坐标系，使 \(B\) 为原点，\(BD\)，\(BA\)，以及垂直于平面 \(ABD\) 的方向分别为 \(x\) 轴、\(y\) 轴、\(z\) 轴. 由已知条件可取
\(B(0,0,0)\)，\(D(1,0,0)\)，\(A(0,1,0)\)，\(C(1,0,1)\)，所以 \(M\left(\frac{1}{2},\frac{1}{2},0\right)\).

平面 \(MBC\) 内的两个不共线向量可取 \(\overrightarrow{BC}=(1,0,1)\) 和 \(\overrightarrow{BM}=\left(\frac{1}{2},\frac{1}{2},0\right)\)，故其一个法向量为
\(\bs{n}=\overrightarrow{BC}\times\overrightarrow{BM}=\left(-\frac{1}{2},\frac{1}{2},\frac{1}{2}\right)\)，可取 \(\bs{n}=(-1,1,1)\). 直线 \(AD\) 的一个方向向量为 \(\overrightarrow{AD}=(1,-1,0)\).

设直线 \(AD\) 与平面 \(MBC\) 所成的角为 \(\theta\)，则
\(\sin\theta=\frac{\left|\bs{n}\cdot\overrightarrow{AD}\right|}{\left|\bs{n}\right|\left|\overrightarrow{AD}\right|}=\frac{\left|(-1)\cdot1+1\cdot(-1)+1\cdot0\right|}{\sqrt{3}\sqrt{2}}=\frac{\sqrt{6}}{3}\).
\end{enumerate}
\end{solution}

\begin{problem}
为回馈顾客，某商场拟通过摸球兑奖的方式对 \(1000\) 位顾客进行奖励，规定：每位顾客从一个装有 \(4\) 个标有面值的球的袋中一次性随机摸出 \(2\) 个球，球上所标的面值之和为该顾客所获的奖励额.

\begin{enumerate}
\item 若袋中所装的 \(4\) 个球中有 \(1\) 个所标的面值为 \(50\text{ 元}\)，其余 \(3\) 个均为 \(10\text{ 元}\)，求：

\begin{enumerate}
\item 顾客所获的奖励额为 \(60\text{ 元}\) 的概率；
\item 顾客所获的奖励额的分布列及数学期望.
\end{enumerate}

\item 商场对奖励总额的预算是 \(60000\text{ 元}\)，并规定袋中的 4 个球只能由标有面值 \(10\text{ 元}\) 和 \(50\text{ 元}\) 的两种球组成，或标有面值 \(20\text{ 元}\) 和 \(40\text{ 元}\) 的两种球组成.为了使顾客得到的奖励总额尽可能符合商场的预算且每位顾客所获的奖励额相对均衡，请对袋中的 4 个球的面值给出一个合适的设计，并说明理由.
\end{enumerate}
\end{problem}

\begin{solution}
设顾客获得的奖励额为随机变量 \(X\).

\begin{enumerate}
\item
\begin{enumerate}
\item 从 \(4\) 个球中一次性取出 \(2\) 个球，共有 \(\mathrm C_4^2=6\) 种等可能取法. 取到标有 \(50\text{ 元}\) 的球和一个标有 \(10\text{ 元}\) 的球时，奖励额为 \(60\text{ 元}\)，这样的取法有 \(\mathrm C_3^1=3\) 种. 因此
\(P(X=60)=\frac{\mathrm C_3^1}{\mathrm C_4^2}=\frac{1}{2}\).

\item 取到两个标有 \(10\text{ 元}\) 的球时，奖励额为 \(20\text{ 元}\)，有 \(\mathrm C_3^2=3\) 种取法；取到标有 \(50\text{ 元}\) 的球和一个标有 \(10\text{ 元}\) 的球时，奖励额为 \(60\text{ 元}\)，有 \(3\) 种取法. 所以 \(X\) 的分布列为
\begin{center}
\begin{tabular}{c|cc}
\(X\)&\(20\)&\(60\)\\
\hline
\(P\)&\(\frac{1}{2}\)&\(\frac{1}{2}\)
\end{tabular}
\end{center}
数学期望为 \(E(X)=20\times\frac{1}{2}+60\times\frac{1}{2}=40\text{ 元}\).
\end{enumerate}

\item 设四个球的面值为 \(v_1,v_2,v_3,v_4\). 每个球恰好出现在 \(\mathrm C_3^1=3\) 种取法中，因此随机取出两个球的奖励额的数学期望为
\(E(X)=\frac{3(v_1+v_2+v_3+v_4)}{\mathrm C_4^2}=\frac{v_1+v_2+v_3+v_4}{2}\).

要使 \(1000E(X)=60000\)，需要 \(E(X)=60\text{ 元}\)，即四个球的面值之和为 \(120\text{ 元}\). 若面值由 \(10\text{ 元}\) 和 \(50\text{ 元}\) 组成，设标有 \(50\text{ 元}\) 的球有 \(m\) 个，则
\(10(4-m)+50m=120\)，解得 \(m=2\). 因此四个球应设计为 \((10,10,50,50)\)，此时 \(X\) 取 \(20\text{ 元},60\text{ 元},100\text{ 元}\)，对应概率分别为 \(\frac{1}{6},\frac{2}{3},\frac{1}{6}\)，且
\(D(X)=\frac{1}{6}(20-60)^2+\frac{2}{3}(60-60)^2+\frac{1}{6}(100-60)^2=\frac{1600}{3}\).

若面值由 \(20\text{ 元}\) 和 \(40\text{ 元}\) 组成，设标有 \(40\text{ 元}\) 的球有 \(m\) 个，则
\(20(4-m)+40m=120\)，解得 \(m=2\). 因此四个球应设计为 \((20,20,40,40)\)，此时 \(X\) 取 \(40\text{ 元},60\text{ 元},80\text{ 元}\)，对应概率分别为 \(\frac{1}{6},\frac{2}{3},\frac{1}{6}\)，且
\(D(X)=\frac{1}{6}(40-60)^2+\frac{2}{3}(60-60)^2+\frac{1}{6}(80-60)^2=\frac{400}{3}\).

两种设计的数学期望都为 \(60\text{ 元}\)，均使奖励总额的期望为 \(60000\text{ 元}\)，但 \(\frac{400}{3}<\frac{1600}{3}\)，所以选择 \((20,20,40,40)\)，此时顾客获得的奖励额相对更均衡.
\end{enumerate}
\end{solution}

\begin{problem}
已知双曲线 \(E: \frac{x^2}{a^2} - \frac{y^2}{b^2} = 1\)（\(a > 0\)，\(b > 0\)）的两条渐近线分别为 \(l_1: y = 2x\)，\(l_2: y = -2x\).

\begin{enumerate}
\item 求双曲线 \(E\) 的离心率；
\item 如图，\(O\) 为坐标原点，动直线 \(l\) 分别交直线 \(l_1\)，\(l_2\) 于 \(A\)，\(B\) 两点（\(A\)，\(B\) 分别在第一，四象限），且 \(\triangle OAB\) 的面积恒为 \(8\)，试探究：是否存在总与直线 \(l\) 有且只有一个公共点的双曲线 \(E\)？若存在，求出双曲线 \(E\) 的方程；若不存在，说明理由.
\end{enumerate}

\bitmapfigure[max width=0.35\linewidth]{img/2014/fujian_science/q19_fig1.png}
\end{problem}

\begin{solution}
\begin{enumerate}
\item 双曲线的渐近线方程为 \(y=\pm\frac{b}{a}x\)，由题设 \(\frac{b}{a}=2\)，得 \(b=2a\). 设双曲线的半焦距为 \(c_0\)，则 \(c_0^2=a^2+b^2=5a^2\)，所以离心率为 \(\frac{c_0}{a}=\sqrt{5}\).

\item 设 \(A(t,2t)\)，\(B(s,-2s)\)，其中 \(t>0\)，\(s>0\). 由三角形面积条件，
\(8=S_{\triangle OAB}=\frac{1}{2}\left|\begin{vmatrix}t&2t\\s&-2s\end{vmatrix}\right|=2ts\)，所以 \(ts=4\)，即 \(s=\frac{4}{t}\).

直线 \(l\) 过 \(A\)，\(B\)，其方程为 \(2(t+s)x+(s-t)y=4ts\). 代入 \(s=\frac{4}{t}\)，并整理得
\(\frac{t^2+4}{8t}x+\frac{4-t^2}{16t}y=1\).

考虑双曲线 \(E_0:\frac{x^2}{4}-\frac{y^2}{16}=1\). 对任意 \(t>0\)，令
\(P_t\left(\frac{t^2+4}{2t},\frac{t^2-4}{t}\right)\). 由于
\(\frac{1}{4}\left(\frac{t^2+4}{2t}\right)^2-\frac{1}{16}\left(\frac{t^2-4}{t}\right)^2=1\)，所以 \(P_t\) 在 \(E_0\) 上.

双曲线 \(E_0\) 在点 \(P_t(x_t,y_t)\) 处的切线方程为 \(\frac{x_tx}{4}-\frac{y_ty}{16}=1\). 代入 \(x_t=\frac{t^2+4}{2t}\)，\(y_t=\frac{t^2-4}{t}\)，该切线正好化为
\(\frac{t^2+4}{8t}x+\frac{4-t^2}{16t}y=1\)，即为直线 \(l\). 因而每一条满足条件的直线 \(l\) 都与 \(E_0\) 相切，只有一个公共点.

所以存在这样的双曲线，其方程为 \(\frac{x^2}{4}-\frac{y^2}{16}=1\).
\end{enumerate}
\end{solution}

\begin{problem}
已知函数 \( f(x) = \e^{x} - ax \)（\(a\) 为常数）的图象与 \(y\) 轴交于点 \(A\)，曲线 \(y = f(x)\) 在点 \(A\) 处的切线斜率为 \(-1\).

\begin{enumerate}
\item 求 \(a\) 的值及函数 \(f(x)\) 的极值；
\item 证明：当 \(x > 0\) 时，\(x^2 < \e^{x}\)；
\item 证明：对任意给定的正数 \(c\)，总存在 \(x_0\)，使得当 \(x \in (x_0, +\infty)\) 时，恒有 \(x^2 < c\e^{x}\).
\end{enumerate}
\end{problem}

\begin{solution}
\begin{enumerate}
\item 因为曲线与 \(y\) 轴交于点 \(A\)，所以 \(A=(0,f(0))=(0,1)\). 又 \(f'(x)=\e^x-a\)，故点 \(A\) 处的切线斜率为 \(f'(0)=1-a=-1\)，解得 \(a=2\).

此时 \(f(x)=\e^x-2x\)，且 \(f'(x)=\e^x-2\). 当 \(x<\ln2\) 时 \(f'(x)<0\)，当 \(x>\ln2\) 时 \(f'(x)>0\)，所以 \(f(x)\) 在 \(x=\ln2\) 处取得极小值，极小值为 \(f(\ln2)=2-2\ln2\).

\item 令 \(g(x)=\e^x-x^2\)，则 \(g'(x)=\e^x-2x=f(x)\). 由第（1）问，\(f(x)\geqslant2-2\ln2\). 又因为 \(2<\e\)，所以 \(\ln2<1\)，从而 \(2-2\ln2>0\)，故 \(g'(x)>0\). 又 \(g(0)=1\)，故当 \(x>0\) 时 \(g(x)>g(0)=1>0\)，即 \(x^2<\e^x\).

\item 取 \(x_0=\frac{16}{c}\). 当 \(c\geqslant1\) 时，由第（2）问可得当 \(x>x_0>0\) 时 \(x^2<\e^x\leqslant c\e^x\).

当 \(0<c<1\) 且 \(x>x_0\) 时，由第（2）问应用于 \(\frac{x}{2}>0\)，得 \(\frac{x^2}{4}<\e^{x/2}\). 又因为 \(x>\frac{16}{c}\)，所以 \(\frac{x^2}{4}>\frac{64}{c^2}>\frac{4}{c}\)，从而 \(\e^{x/2}>\frac{4}{c}\). 因此 \(4\e^{x/2}<c\e^x\)，结合 \(x^2<4\e^{x/2}\)，得到 \(x^2<c\e^x\).

综上，对任意正数 \(c\)，当 \(x\in\left(\frac{16}{c},+\infty\right)\) 时恒有 \(x^2<c\e^x\).
\end{enumerate}
\end{solution}

\section{选考题：解答题}

\begin{problem}
已知矩阵 \(A\) 的逆矩阵 \(A^{-1} = \begin{pmatrix} 2 & 1 \\ 1 & 2 \end{pmatrix}\).

\begin{enumerate}
\item 求矩阵 \(A\)；
\item 求矩阵 \(A^{-1}\) 的特征值以及属于每个特征值的一个特征向量.
\end{enumerate}
\end{problem}

\begin{solution}
\begin{enumerate}
\item 由
\(A^{-1}=\begin{pmatrix}2&1\\1&2\end{pmatrix}\)，且 \(\det A^{-1}=4-1=3\)，可得
\(A=(A^{-1})^{-1}=\frac{1}{3}\begin{pmatrix}2&-1\\-1&2\end{pmatrix}\).

\item 设 \(A^{-1}\) 的特征值为 \(\lambda\)，则
\(\det\left(\begin{pmatrix}2&1\\1&2\end{pmatrix}-\lambda I\right)=(2-\lambda)^2-1=(\lambda-1)(\lambda-3)=0\).
所以特征值为 \(1\) 和 \(3\).

当 \(\lambda=1\) 时，特征方程为 \(x+y=0\)，可取特征向量 \(\begin{pmatrix}1\\-1\end{pmatrix}\)；当 \(\lambda=3\) 时，特征方程为 \(-x+y=0\)，可取特征向量 \(\begin{pmatrix}1\\1\end{pmatrix}\).
\end{enumerate}
\end{solution}

\begin{problem}
已知直线 \(l\) 的参数方程为 \(\begin{cases}
x=a-2t,\\
y=-4t,
\end{cases}\)（\(t\) 为参数），圆 \(C\) 的参数方程为 \(\begin{cases}
x=4\cos\theta,\\
y=4\sin\theta,
\end{cases}\)（\(\theta\) 为参数）.

\begin{enumerate}
\item 求直线 \(l\) 和圆 \(C\) 的普通方程；
\item 若直线 \(l\) 与圆 \(C\) 有公共点，求实数 \(a\) 的取值范围.
\end{enumerate}
\end{problem}

\begin{solution}
\begin{enumerate}
\item 由直线 \(l\) 的参数方程 \(y=-4t\)，得 \(t=-\frac{y}{4}\). 代入 \(x=a-2t\)，得 \(x=a+\frac{y}{2}\)，所以直线 \(l\) 的普通方程为 \(2x-y-2a=0\).

由圆 \(C\) 的参数方程，消去参数 \(\theta\) 得 \(x^2+y^2=16\)，这就是圆 \(C\) 的普通方程.

\item 直线 \(l\) 与圆 \(C\) 有公共点，当且仅当原点到直线 \(l\) 的距离不大于圆的半径 \(4\). 因此
\(\frac{\left|-2a\right|}{\sqrt{2^2+(-1)^2}}\leqslant4\)，即 \(\frac{2|a|}{\sqrt5}\leqslant4\). 解得 \(|a|\leqslant2\sqrt5\)，所以 \(-2\sqrt5\leqslant a\leqslant2\sqrt5\).
\end{enumerate}
\end{solution}

\begin{problem}
已知定义在 \(\R\) 上的函数 \(f(x) = |x + 1| + |x - 2|\) 的最小值为 \(a\).

\begin{enumerate}
\item 求 \(a\) 的值；
\item 若 \(p\)，\(q\)，\(r\) 为正实数，且 \(p + q + r = a\)，求证：\(p^2 + q^2 + r^2 \geqslant 3\).
\end{enumerate}
\end{problem}

\begin{solution}
\begin{enumerate}
\item 分三种情况讨论：当 \(x<-1\) 时，\(f(x)=-2x+1>3\)；当 \(-1\leqslant x\leqslant2\) 时，\(f(x)=x+1+2-x=3\)；当 \(x>2\) 时，\(f(x)=2x-1>3\). 因此 \(f(x)\) 的最小值为 \(3\)，即 \(a=3\).

\item 由 \(p+q+r=a=3\)，以及柯西不等式，得
\(p^2+q^2+r^2\geqslant\frac{(p+q+r)^2}{1^2+1^2+1^2}=\frac{9}{3}=3\).
\end{enumerate}
\end{solution}
